\section{Experiments with the general algorithms}\label{sec:experiments}

\paragraph{Implementation.}
All algorithms are implemented in C++17 in a single header
(\code{convex/kdisp.hpp}); geometric decisions use only comparisons of $64$-bit
squared distances and the $128$-bit in-circle determinant. Besides
\cref{alg:capped,alg:certified} the code contains two baselines: the same
dynamic program without the cap ($K=n-2$, cf.\ \cref{rem:cap}), and the direct
\emph{triple-state} program that stores a value for every chord and apex and
scans all sub-apexes, which takes $O(n^4)$ in-circle tests per decision. All
optimisation variants use the search window unless stated otherwise.
Experiments were run on one core of a virtual machine (x86-64, GCC~13.3,
\code{-O2}); times are medians of three runs for $n\le200$ and single runs
for larger instances, and exclude reading the input and building the distance
matrix.

\paragraph{Instances.}
Every instance is generated by sampling a cloud of random integer points,
taking its strictly convex hull, and keeping $n$ random hull vertices. The
distributions are: uniform in a square and Gaussian (small hulls only),
uniform in a disk, uniform in a thin annulus, uniform in a $1{:}20$ ellipse
(\code{thin}), rounded points on a circle, exactly cocircular lattice points
(the $972$ lattice points on a circle of radius $1185665$, a worst case for
in-circle tests), and a thin lens bounded by two flat parabolic arcs
(\code{twochain}), the shape on which greedy rules fail. Coordinates are
bounded by $1.2\cdot10^{6}$. There are four splits: \emph{verify}
($6\le n\le14$, $\nVerify$ instances), \emph{medium} ($20\le n\le40$,
$\nMedium$ instances), \emph{scaling} ($25\le n\le400$, $k\in\{4,8,16,32\}$,
$\nScaling$ instances), and a sweep over $k$ at $n=200$.

\paragraph{Correctness.}
On every instance of every split, the capped program with and without the
window, the uncapped program, the certified algorithm and (for $n\le100$) the
triple-state program returned the same optimum, and every reported set was
checked to be a $k$-set with exactly that minimum distance
(\cref{tab:correct}). On \emph{verify} the optimum also agrees with a
branch-and-bound over all $k$-subsets, with an independent brute force
written in Python ($\nVerifyPy/\nVerify$), and with an independent
general-dimension solver based on clique enumeration and matrix products that
is part of the same code base ($\nVerifyFmm/\nVerify$). On \emph{medium},
exhaustive search was run for $k\le5$ ($\nMediumBrute$ instances). The
explicit instances of \cref{sec:greedy} are re-verified by a separate script
that shares no code with the solvers. In addition, a randomised self-test
compares all solvers and \cref{cor:mis} with exhaustive search on tens of
thousands of small instances, including degenerate ones with many equal
distances and cocircular quadruples.

\begin{table}[tb]
\centering
\caption{Cross-validation. ``agree'': all solvers return the same optimum and
valid optimal sets; ``exhaustive'': number of instances on which this value
was also confirmed by exhaustive search.}
\label{tab:correct}
\small
\begin{tabular}{@{}llrrrr@{}}
\toprule
split & distribution & instances & $n$ & agree & exhaustive\\
\midrule
\tabinput{generated/table_correctness}
\bottomrule
\end{tabular}
\end{table}

\paragraph{Running time.}
\Cref{fig:time} shows the median running times on the
\emph{scaling} split. The triple-state baseline grows roughly like $n^4$ and is
already $\tripleSpeedupHundred$ times slower than the capped program at
$n=100$.
At the optimum, a median of $\farTripleMedian\%$ of all triples are
$r$-far (between $\farTripleMin\%$ and $\farTripleMax\%$), so the term
$\Tr\log k$ dominates in practice.

\begin{figure}[tb]
\centering
\begin{tikzpicture}
\begin{groupplot}[group style={group size=4 by 1, horizontal sep=0.75cm},
  width=0.285\textwidth, height=0.36\textwidth, xmode=log, ymode=log,
  xlabel={$n$}, xtick={25,50,100,200,400}, xticklabels={25,50,100,200,400}, x tick label style={font=\tiny},
  tick label style={font=\scriptsize}, label style={font=\small},
  title style={font=\small}, grid=major, grid style={gray!25},
  legend style={font=\scriptsize, at={(2.55,1.22)}, anchor=south, legend columns=4, draw=gray!50, /tikz/every even column/.append style={column sep=8pt}},
  every axis plot/.append style={mark size=1.5pt}]
\nextgroupplot[title={$k=4$}, ylabel={time [ms]}]
  \addplot[black, mark=triangle*] table {generated/plot_k4_triple.dat};
  \addplot[orange!85!black, mark=diamond*] table {generated/plot_k4_uncapped.dat};
  \addplot[blue, mark=*] table {generated/plot_k4_capped.dat};
  \addplot[red!80!black, mark=square*] table {generated/plot_k4_certified.dat};
  \legend{triple-state, uncapped, capped, certified}
\nextgroupplot[title={$k=8$}]
  \addplot[black, mark=triangle*] table {generated/plot_k8_triple.dat};
  \addplot[orange!85!black, mark=diamond*] table {generated/plot_k8_uncapped.dat};
  \addplot[blue, mark=*] table {generated/plot_k8_capped.dat};
  \addplot[red!80!black, mark=square*] table {generated/plot_k8_certified.dat};
\nextgroupplot[title={$k=16$}]
  \addplot[black, mark=triangle*] table {generated/plot_k16_triple.dat};
  \addplot[orange!85!black, mark=diamond*] table {generated/plot_k16_uncapped.dat};
  \addplot[blue, mark=*] table {generated/plot_k16_capped.dat};
  \addplot[red!80!black, mark=square*] table {generated/plot_k16_certified.dat};
\nextgroupplot[title={$k=32$}]
  \addplot[black, mark=triangle*] table {generated/plot_k32_triple.dat};
  \addplot[orange!85!black, mark=diamond*] table {generated/plot_k32_uncapped.dat};
  \addplot[blue, mark=*] table {generated/plot_k32_capped.dat};
  \addplot[red!80!black, mark=square*] table {generated/plot_k32_certified.dat};
\end{groupplot}
\end{tikzpicture}
\caption{Median time to solve $k$-dispersion exactly, over all distributions
of the \emph{scaling} split (log--log). The triple-state baseline was run for
$n\le100$ and the uncapped program for $n\le200$.}
\label{fig:time}
\end{figure}

\paragraph{Effect of the cap.}
\Cref{tab:ksweep} compares the capped and the uncapped program at $n=200$.

\begin{table}[tb]
\centering
\caption{Sweep over $k$ at $n=200$ (median over six distributions): time [ms]
of the uncapped and the capped program and of the certified algorithm, and
in-circle tests per exact decision of the capped program (in millions).}
\label{tab:ksweep}
\small
\begin{tabular}{@{}rrrrrr@{}}
\toprule
$k$ & uncapped & capped & uncapped/capped & certified & in-circle tests [$10^6$]\\
\midrule
\tabinput{generated/table_ksweep}
\bottomrule
\end{tabular}
\end{table}

\paragraph{Certification.}
\Cref{tab:certified} reports how often the heuristic phase of
\cref{alg:certified} already attains the optimum, over all
$\certTotal$ instances of the verify, medium and scaling splits. This happens
on $\certYes$ instances, and on every instance drawn from the square,
Gaussian, disk, annulus, circle and cocircular distributions. In these cases
a single exact decision is made, and the certified algorithm is
$\certSpeedup$ times faster than the capped binary search in the median
(interquartile range $\certSpeedupLo$--$\certSpeedupHi$). The farthest-first
set alone is optimal far less often, so the seeded greedies do most of the
work. The output is exact in every case.

\begin{table}[tb]
\centering
\caption{Quality of the heuristic lower bounds. ``FF optimal'': the
farthest-first set attains $\rstar$; ``certified'': the heuristic phase of
\cref{alg:certified} attains $\rstar$, so that one exact decision suffices.
The last two columns give the mean and the maximum relative gap
$1-g/\rstar$ of the farthest-first value $g$.}
\label{tab:certified}
\small
\begin{tabular}{@{}lrrrrrr@{}}
\toprule
distribution & instances & FF optimal & certified & rate & mean gap & max gap\\
\midrule
\tabinput{generated/table_certified}
\bottomrule
\end{tabular}
\end{table}


\section{Experiments for the arc property}\label{app:arcexp}

\paragraph{Verification.}
The implementation (\code{arc/arcdisp.hpp}) was validated before any timing.
On $\arcSmallInst$ instances with $n\le16$ from ten families, at every
candidate radius $r$ (and at $r^2+1$) at which the arc property holds, three
implementations of the decision procedure agree with exhaustive search for
all $k=2,\dots,n$ ($\arcDchecks$ comparisons), and \textsc{O1}, \textsc{O2}
and the adaptive variant return the exhaustive optimum with a valid witness
($\arcOptChecks$ instance--$k$ pairs). On $\arcBigInst$ instances with
$20\le n\le150$ they agree with \cref{alg:capped} ($\arcBigChecks$ checks).
An independent Python script recomputes the optima of exported instances by
brute force and checks \cref{thm:arc} directly ($\arcPyAtwo$ checks). No
discrepancy occurred.

\paragraph{Running time.}
Unimodal instances are points at jittered equally spaced angles on a circle,
rounded to integers; every instance is verified to be strictly convex and
unimodal by the exact $O(n\log n)$ test. With coordinates bounded by $2^{29}$
such instances exist for $n\le50000$; for $n\ge10^5$ we use a circle of radius
$2^{60}$ and $128$-bit squared distances. \textsc{O1} needs $\Theta(n^2)$
memory and was run for $n\le10^4$. Every reported witness was
checked by an exact closest-pair computation. \Cref{tab:arcscaling} shows the
medians over three instances. \textsc{O2} solves instances with a million
points in about $\arcMillion$ seconds, and its running time grows like
$n^{\arcSlope}$ between $n=10^3$ and $n=10^6$, compared with
$n^{\arcSlopeOone}$ for \textsc{O1}. At $n=400$ it is between $300$ ($k=4$) and $40\,000$ ($k=64$) times faster
than the cubic dynamic program.

\begin{table}[tb]
\centering
\caption{Median running time [ms] on unimodal instances (three instances per
cell). \textsc{DP}: \cref{alg:capped} with binary search. ``--'': not run.}
\label{tab:arcscaling}
\small
\begin{tabular}{@{}rr|rrr|r@{}}
\toprule
 & & \multicolumn{3}{c|}{$64$-bit distances} & $128$-bit\\
$n$ & $k$ & \textsc{DP} & \textsc{O1} & \textsc{O2} & \textsc{O2}\\
\midrule
\tabinput{generated/arc_scaling}
\bottomrule
\end{tabular}
\end{table}



